%! The Idea of Elimination — Unit 2, Lesson 2.1 — Guided Notes
\documentclass[10pt]{article}
\usepackage{linalg-article}
\usepackage{linalg-boxes}
\newcommand{\TallMath}[1]{$\displaystyle #1\rule[-1.4em]{0pt}{3.2em}$}
\newcommand{\xx}{\mathbf{x}}
\newcommand{\bb}{\mathbf{b}}

\begin{document}

\pageheader{Unit 2, Lesson 2.1}{Guided Notes}

\vspace{0.12in}

\begin{objectivebox}
\textbf{By the end of this lesson, I will be able to\dots}
\begin{itemize}\setlength{\itemsep}{2pt}
  \item find the \emph{pivot} and the \emph{multiplier} $\ell=$ entry $\div$ pivot, and use them to make a zero;
  \item reach \emph{upper-triangular} form, then \emph{back-substitute} and check by rebuilding $\bb$;
  \item fix a \emph{zero pivot} with a \emph{row exchange}, and read a breakdown as no or infinitely many solutions.
\end{itemize}
\end{objectivebox}

\vspace{0.10in}

\begin{vocabbox}
\small
Fill in each term as we build it together.
\par\vspace{2pt}
\termblanklong{Pivot}
\termblanklong{Multiplier $\ell$}
\termblanklong{Upper-triangular system}
\termblanklong{Back-substitution}
\termblanklong{Row exchange}
\end{vocabbox}

\begin{hookbox}
Two systems, side by side.
\[
  \text{(Lesson 2.0)}\quad
  \begin{array}{r@{}c@{}l} x &{}+ 2y ={}& 5 \\ 3x &{}+ 2y ={}& 11 \end{array}
  \qquad\qquad
  \text{(today)}\quad
  \begin{array}{r@{}c@{}l} 2x &{}+ 3y ={}& 13 \\ 4x &{}+ 7y ={}& 27 \end{array}
\]
\begin{itemize}\setlength{\itemsep}{1pt}
  \item In the \emph{left} system, one subtraction removes $y$. Why does it work so easily? \writeline
  \item In the \emph{right} system, the $x$-terms are $2$ and $4$. What could you do to the first
        equation \emph{before} subtracting so that $x$ disappears? \writeline
\end{itemize}
\end{hookbox}

\begin{notesbox}{1. Pivots and Multipliers}
When the coefficients don't already match, we \textbf{make} them match. Work with the right-hand
system above: $2x+3y=13$ and $4x+7y=27$.

\vspace{3pt}
The \textbf{pivot} is the first nonzero coefficient of the top row --- the coefficient we build on.
Here the pivot is \blank{1.0cm} (in front of $x$).

\vspace{3pt}
To clear the $x$-term below the pivot, we need to know \emph{how many pivot rows} to subtract. That
number is the \textbf{multiplier}:
\[
  \ell = \frac{\text{entry to eliminate}}{\text{pivot}} = \frac{4}{2} = \blank{1.0cm}.
\]
Now do the \textbf{elimination step} --- replace row~2 by (row~2) $-\;\ell\times$(row~1):
\[
  \begin{array}{r@{}c@{}l}
    (4x + 7y) - 2\,(2x + 3y) &{}={}& 27 - 2\,(13) \\[2pt]
    (4-4)x + (7-6)y &{}={}& 27 - 26 \\[2pt]
    y &{}={}& \blank{1.0cm}
  \end{array}
\]
The $x$-term vanished \emph{by design}: that is exactly what the multiplier was chosen to do.
\end{notesbox}

\begin{notesbox}{2. Triangular Form, Then Back-Substitute}
After that one step the system looks like a staircase --- all zeros below the pivot. This is
\textbf{upper-triangular} form:
\[
  \begin{array}{r@{}c@{}l}
    2x &{}+ 3y ={}& 13 \\
       & y ={}& 1
  \end{array}
\]
The bottom equation already gives $y$. \textbf{Back-substitute} it into the top row and solve
\emph{upward}:
\[
  2x + 3(1) = 13 \;\Longrightarrow\; 2x = \blank{1.0cm} \;\Longrightarrow\; x = \blank{1.0cm}.
\]

\vspace{2pt}
\textbf{Check by rebuilding $\bb$.} The columns of $A$ are $\begin{bsmallmatrix} 2\\4 \end{bsmallmatrix}$
and $\begin{bsmallmatrix} 3\\7 \end{bsmallmatrix}$; scale by the weights and add:
\[
  5\begin{bmatrix} 2 \\ 4 \end{bmatrix} + 1\begin{bmatrix} 3 \\ 7 \end{bmatrix}
  = \begin{bmatrix}\rule{0pt}{1.1em}\blank{1.1cm}\\[2pt]\blank{1.1cm}\end{bmatrix}.
\]
So the solution is $x=\blank{0.9cm}$, $y=\blank{0.9cm}$.
\end{notesbox}

\begin{notesbox}{3. Bigger Systems: The Same Staircase}
Elimination scales up. For three equations, use row~1's pivot to clear the \emph{whole} first
column, then row~2's new pivot to clear below it.
\[
  \begin{array}{r@{}c@{}l}
    x &{}+ y + z ={}& 6 \\
    2x &{}+ 3y + z ={}& 11 \\
    x &{}+ 3y + 6z ={}& 25
  \end{array}
\]
\textbf{Round 1} (pivot $=1$ in row~1). Clear the $x$ from rows~2 and~3:
\begin{itemize}\setlength{\itemsep}{1pt}
  \item row~2 $-\;\ell_{21}\cdot$row~1 with $\ell_{21}=2\div1=\blank{0.8cm}$ \;$\Rightarrow$\; $y - z = -1$;
  \item row~3 $-\;\ell_{31}\cdot$row~1 with $\ell_{31}=1\div1=\blank{0.8cm}$ \;$\Rightarrow$\; $2y + 5z = 19$.
\end{itemize}
\textbf{Round 2} (new pivot $=1$ in row~2). Clear the $y$ from row~3 with
$\ell_{32}=2\div1=\blank{0.8cm}$: \;$(2y+5z)-2(y-z)=19-2(-1)\Rightarrow 7z=21$. The triangle:
\[
  \begin{array}{r@{}c@{}l}
    x &{}+ y + z ={}& 6 \\
     & y - z ={}& -1 \\
     & 7z ={}& 21
  \end{array}
\]
Back-substitute upward: $z=\blank{0.8cm}$, then $y=\blank{0.8cm}$, then $x=\blank{0.8cm}$.
Solution $(x,y,z)=(\,\blank{1.6cm}\,)$.
\end{notesbox}

\boxguard
\begin{notesbox}{4. When a Pivot Is Zero}
Elimination needs a \emph{nonzero} pivot. If the pivot slot holds a $0$, look \emph{below} for a row
with a nonzero in that spot and \textbf{exchange rows}.
\[
  \begin{array}{r@{}c@{}l}
    0x &{}+ 2y ={}& 4 \\ 3x &{}+ y ={}& 5
  \end{array}
  \qquad\xrightarrow{\text{swap the rows}}\qquad
  \begin{array}{r@{}c@{}l}
    3x &{}+ y ={}& 5 \\ & 2y ={}& 4
  \end{array}
\]
After the swap it is already triangular: $2y=4$ gives $y=\blank{0.8cm}$, then $3x+2=5$ gives
$x=\blank{0.8cm}$.

\vspace{3pt}
What if \emph{no} row can be swapped in? Then elimination stalls, and the last line tells the story:
$0 = (\text{a nonzero})$ means \blank{3.2cm}, while $0=0$ means \blank{3.6cm}. (Row exchanges get
their own lesson in 2.4.)
\end{notesbox}

\boxguard
\begin{practicebox}
Solve $A\xx=\bb$ with $A=\begin{bmatrix} 3 & 2 \\ 6 & 7 \end{bmatrix}$ and
$\bb=\begin{bmatrix} 8 \\ 19 \end{bmatrix}$ by elimination.
\begin{enumerate}[leftmargin=1.6em]
  \item Name the pivot and the multiplier $\ell$ that clears the entry below it. \hfill
        pivot $=\blank{1.2cm}$\quad $\ell=\blank{1.2cm}$
  \item Do the elimination step to reach upper-triangular form, and find $y$.
\begin{work}
  (6x+7y)-2(3x+2y) &= 19-16 \\
               3y &= 3 \\
                y &= 1
\end{work}
  \item Back-substitute to find $x$, then \textbf{check} by rebuilding $\bb$.
\begin{work}
  3x+2 &= 8 \\
     x &= 2 \\
  2\begin{bsmallmatrix} 3\\6 \end{bsmallmatrix} + 1\begin{bsmallmatrix} 2\\7 \end{bsmallmatrix}
       &= \begin{bsmallmatrix} 8\\19 \end{bsmallmatrix} \ \ \text{\checkmark}
\end{work}
\end{enumerate}
\end{practicebox}

\end{document}
